%% ICAIF '26 submission — ACM sigconf, double-blind.
%% Compile on Overleaf (acmart is preinstalled) or any TeX Live >= 2020.
%% 8-page limit INCLUDING references. Deadline: Aug 2, 2026 AoE. Via CMT.
%%
%% PENDING numbers are marked \pending{...} and render orange. Search for
%% "pending" before submission — the build is not submittable while any remain.

\documentclass[sigconf,anonymous,review]{acmart}

\usepackage{booktabs}
\usepackage{xcolor}

%% Placeholder macro for numbers awaiting cluster jobs.
\newcommand{\pending}[1]{\textcolor{orange}{\textbf{[PENDING: #1]}}}

%% ACM rights boilerplate — leave as-is for review; CMT/ACM will supply
%% the real block on acceptance.
\setcopyright{none}
\acmConference[ICAIF '26]{7th ACM International Conference on AI in Finance}{November 2026}{TBD}
\acmYear{2026}
\settopmatter{printacmref=false}

\begin{document}

\title{Hard Switching, Soft Gating, or Learned Routing?\\
How to Inject a Market-Regime Signal into a Reinforcement-Learning Portfolio Policy}

%% Anonymous for review — acmart's [anonymous] option suppresses these,
%% but keep placeholders so the camera-ready diff is one commit.
\author{Anonymous Author(s)}
\affiliation{%
  \institution{Anonymous Institution}
  \city{}
  \country{}
}

\begin{abstract}
Reinforcement-learning (RL) agents for portfolio management are typically
trained against a single stationary reward landscape, yet markets alternate
between structurally distinct regimes in which the optimal allocation policy
differs. Given an \emph{external} regime signal --- here, the posterior of a
4-state Gaussian Hidden Markov Model (HMM) fitted on macroeconomic ETFs,
disjoint from the trading universe --- we ask the architectural question:
how should the signal enter the policy network? We compare (i) \emph{hard
switching}, which discretely selects one of four regime-specialist heads;
(ii) \emph{soft gating}, which blends all four heads by the HMM posterior;
and (iii) a \emph{learned router} trained end-to-end without consuming the
posterior directly, against a regime-agnostic baseline, on 38 U.S. equities
(2015--2024). We evaluate under two on-policy algorithms (A2C and PPO), with
seed-paired statistics, an expanding 8-window walk-forward protocol, and
Jobson--Korkie--Memmel tests on daily returns. Three findings survive every
robustness pass. First, the learned router is the worst and least stable
variant under \emph{both} algorithms, collapsing on a subset of seeds ---
evidence that a stable, separately estimated regime signal beats jointly
learned routing. Second, the gate's value is \emph{protocol-dependent}:
trained once and tested across an unseen bear-plus-recovery, soft gating
dominates hard switching in effect size and stability under A2C (Sharpe
0.46 vs.\ 0.14; every soft seed positive while two of three hard seeds
collapse) and ties under PPO's clipped objective (0.774 vs.\ 0.726) ---
but in the walk-forward, where policies are retrained every six months
on expanding data, no gated variant beats the no-gate baseline:
\emph{frequent retraining substitutes for regime awareness}. Third, a
hand-crafted HMM-gated equal-weight rule \emph{destroys} value (Sharpe
0.41 $\rightarrow$ 0.16) --- the signal is not directly monetizable; any
value emerges only through learned exploitation. Finally, a cautionary
methodological result: a sliding-window update scheme had manufactured
an internally consistent, significant pro-gating narrative --- a
drawdown advantage at $p{=}0.0025$, a walk-forward ranking, and a
COVID-crash stress-test story --- all of which reversed or dissolved
under corrected non-overlapping batches. Architectural comparisons in RL
finance can be artifacts of training mechanics; we propose mechanics
ablations as a standard robustness pass.
\end{abstract}

\begin{CCSXML}
<ccs2012>
<concept>
<concept_id>10010405.10003550.10003554</concept_id>
<concept_desc>Applied computing~Portfolio management</concept_desc>
<concept_significance>500</concept_significance>
</concept>
<concept>
<concept_id>10010147.10010257.10010258.10010261</concept_id>
<concept_desc>Computing methodologies~Reinforcement learning</concept_desc>
<concept_significance>500</concept_significance>
</concept>
</ccs2012>
\end{CCSXML}
\ccsdesc[500]{Applied computing~Portfolio management}
\ccsdesc[500]{Computing methodologies~Reinforcement learning}

\keywords{regime-aware reinforcement learning, portfolio management, hidden
Markov models, mixture of experts, soft gating, reproducibility}

\maketitle

%% ===================================================================
\section{Introduction}
%% ===================================================================

Deep reinforcement learning has become a standard alternative to classical
mean--variance optimisation for portfolio allocation, learning sequential,
state-dependent policies directly from data~\cite{jiang2017drl,liu2020finrl}.
A persistent weakness is non-stationarity: markets transition between
qualitatively different regimes --- sustained bulls, sharp crashes, choppy
recoveries, prolonged bears --- each of which may demand a different
allocation policy. A single network trained to perform well \emph{on
average} risks being mediocre in every regime rather than excellent in any.

This motivates \emph{regime-aware RL} (RA-RL), in which the agent receives
an explicit signal about the current regime. Prior work has largely focused
on whether such signals help at all. We address the question one level
down, which prior work leaves open: \textbf{given a regime signal, how
should it be incorporated into the policy network?} We compare three
mechanisms within one framework, holding data, reward, universe, and
training budget fixed:
\begin{enumerate}
  \item \textbf{Hard RA-RL} --- the maximum-a-posteriori regime label
        discretely selects one of $K$ regime-specialist policy heads;
  \item \textbf{Soft RA-RL} --- the full $K$-dimensional HMM posterior
        blends all heads' outputs as a probability-weighted mixture;
  \item \textbf{Learned Router} --- a gating MLP trained end-to-end
        produces its own routing weights, isolating the value of the
        \emph{external} signal from the soft-blending mechanism itself.
\end{enumerate}
All three are benchmarked against a regime-agnostic baseline, classical
allocators (equal-weight, Markowitz, buy-and-hold), and --- critically ---
a \emph{non-RL} rule that consumes the same HMM signal, which tests whether
RL is needed at all.

The regime signal is a 4-state Gaussian HMM~\cite{hamilton1989} fitted,
unsupervised and offline, on six macroeconomic ETFs disjoint from the
trading universe. This separates regime \emph{detection} (stable, external,
estimated once) from regime \emph{exploitation} (the policy's use of the
signal) --- a separation whose value the learned-router ablation directly
quantifies.

Because RL results in finance are notoriously fragile, we subject every
claim to three robustness passes that we believe should be standard:
(a)~\textbf{algorithm robustness} --- every comparison is run under both
A2C~\cite{mnih2016a2c} and PPO~\cite{schulman2017ppo}, which differ exactly
in trust-region control; (b)~\textbf{seed-paired statistics} --- all
variants share seeds, so we use paired $t$ and Wilcoxon signed-rank tests,
plus Jobson--Korkie--Memmel~\cite{jobson1981,memmel2003} tests on daily
returns; and (c)~\textbf{mechanics robustness} --- we re-ran the full
comparison after replacing a stride-1 sliding-window update buffer with
non-overlapping batches, and report which claims survived. One headline
claim did not: soft gating's maximum-drawdown advantage, the most
statistically significant result under the original mechanics
($p{=}0.0025$, paired), reversed sign under the corrected mechanics under
\emph{both} algorithms. We report this openly as a contribution:
architectural conclusions in RL-for-finance can be manufactured by the
update scheme, and mechanics ablations belong in the evaluation protocol.

\textbf{Contributions.}
(1)~The first controlled comparison, within one RL portfolio framework, of
an externally supplied regime gate, a discretised version of the same gate,
and a learned end-to-end router.
(2)~An algorithm-robustness result: the learned router is the worst and
least stable variant under both A2C and PPO, while the soft-vs-hard gap is
algorithm-\emph{dependent} --- large without trust-region control, a tie
with it.
(3)~A ``why RL at all'' baseline: hand-crafted rules on the same HMM signal
destroy value, while the soft-gated agent beats plain equal-weight on the
single-period protocol.
(4)~A reproducibility case study: a training-mechanics artifact that
produced a spurious $p{=}0.0025$ architectural claim, and the protocol that
caught it.

%% ===================================================================
\section{Related Work}
%% ===================================================================

\textbf{RL for portfolio management.}
Policy-gradient and actor-critic methods have been applied to portfolio
allocation~\cite{jiang2017drl,moody2001,mnih2016a2c}, often via
frameworks such as FinRL~\cite{liu2020finrl}, with recent work refining
the reward side of the problem~\cite{karzanov2025regret}; these
typically train a single policy without explicit regime conditioning.
Sun et al.~\cite{sun2022deeptrader} route stock predictions through
uncertainty-aware expert sub-models, but their routing addresses
\emph{model} uncertainty over predictors, not market-regime conditioning
of an RL policy.

\textbf{Regime-switching models.}
Hamilton's Markov-switching framework~\cite{hamilton1989} established that
returns exhibit persistent discrete regimes well captured by latent-state
models; HMM-derived probabilities have since been used for allocation
mostly within classical optimisation, including point-in-time macro-regime
models for factor investing~\cite{nuriyev2024macro} and regime-aware
sample weighting for forecasting~\cite{jang2025regime}. HMM generative
models have also been used to simulate equity dynamics for downstream
decision systems~\cite{istiaque2023hmm}. Closest to us, Nixon
Raj~\cite{nixonraj2025} benchmarks PPO variants under \emph{latent} regime
shifts inferred implicitly, and Verma et al.~\cite{verma2026} pair a
Markov-switching model with RL for three-asset tactical allocation.
Neither isolates our architectural question: given an \emph{explicit},
externally estimated regime posterior, what is the best way to inject it
into the policy? We hold the universe to a single asset class precisely to
keep this comparison clean.

\textbf{Mixture-of-experts and gating.}
MoE~\cite{jacobs1991} and its large-scale
descendants~\cite{shazeer2017moe} train the gate jointly with the experts
on the same objective. Soft RA-RL borrows the specialists-plus-blend
structure but deliberately departs from MoE: the blend weights are the
posterior of an HMM fitted offline on disjoint data and receive no
gradient from the portfolio reward. MoEDRLPM~\cite{moedrlpm2025} is a
contemporaneous RL portfolio model with an expert pool and a
\emph{learned} router in the classical MoE sense --- architecturally our
Learned Router arm, whose instability we document under two algorithms;
learned MoE gating has likewise been applied to trade
execution~\cite{xu2025moe}.

\textbf{Latent context inference in RL.}
An alternative to an external signal is to infer context from interaction,
as in VariBAD~\cite{zintgraf2020varibad}. Our Learned Router is
conceptually in this family, and its seed collapses under both A2C and PPO
are consistent with the known difficulty of learning context inference
jointly with control.

\textbf{The 1/N puzzle.}
DeMiguel et al.~\cite{demiguel2009} showed naive equal weighting often
beats optimised portfolios out of sample. We replicate this in the
walk-forward protocol, scope our claims accordingly, and add a twist: on
the 2022--2024 single-period test, the soft-gated agent beats 1/N.

%% ===================================================================
\section{Methodology}
%% ===================================================================

\subsection{Problem Formulation}
Portfolio allocation is a continuous-control MDP. Each trading day $t$ the
agent observes state $s_t$ and outputs target weights $w_t$ over $N{=}38$
assets with $\sum_i w_{t,i}=1$ and $w_{t,i}\in[-0.05,0.20]$ (limited
shorting, 20\% position cap). Raw policy outputs are mapped to this
bounded simplex by Euclidean projection, computed by bisection on the
projection's Lagrange multiplier.

\subsection{Regime Detection}
A $K{=}4$-state Gaussian HMM is fitted on daily features of six
macroeconomic ETFs (SPY, DBC, LQD, EMB, TLT, TIP --- equities,
commodities, credit, EM debt, duration, inflation-protected), using
training-period data only (no lookahead). States are ordered by a
momentum-minus-volatility score into an interpretable bear-to-bull
ordering. The HMM emits either a hard MAP label or the full posterior
$p_t\in\Delta^3$, depending on the consuming architecture.
$K{=}4$ is selected by BIC: it beats $K{=}5$ (11948 vs.\ 11998), and
larger $K$ wins only via degenerate states with 0.3--2\% occupancy
(Figure~\ref{fig:bic}). \pending{K=2/3 end-to-end ablation if budget
allows; currently queued}

\subsection{Architectural Variants}
All variants share a 1{,}596-dimensional base state: the flattened
$38{\times}38$ return covariance (1{,}444) plus MACD/RSI/CCI/ADX per asset
(152). They differ only in what regime information is appended and how the
actor uses it. All actors share a two-layer 256-unit ReLU MLP trunk and a
state-independent log-std Gaussian head.

\textbf{Baseline (no gate).} Single Gaussian head, no regime input.

\textbf{Hard RA-RL.} The MAP label (scalar) is appended; four specialist
heads share the trunk; the label masks all but one head per step, so only
one head receives gradient per update.

\textbf{Soft RA-RL.} The 4-dim posterior is appended; the same four heads
are all evaluated every step and their outputs combined as
$a_t=\sum_k p_t(k)\, a_t^{(k)}$ --- the posterior \emph{is} the gate; no
gating parameters exist.

\textbf{Learned Router.} Identical experts-and-blend structure, but a
router MLP consuming the trunk features and the raw posterior produces the
blend weights, trained end-to-end with the policy. This isolates the value
of using the external signal \emph{directly}.

\subsection{Reward}
All variants optimise the same mean--variance-penalised reward
\[
r_t = \big(\Delta_t - \tfrac{1}{2}\lambda\, w_t^{\top}\Sigma_t w_t
      - c\cdot \mathrm{TO}_t - \kappa \textstyle\sum_i w_{t,i}^2\big)\times s,
\]
with realised return $\Delta_t$, covariance $\Sigma_t$, turnover penalty
$c{=}10^{-4}$, concentration penalty $\kappa{=}0.005$, risk aversion
$\lambda{=}0.5$, and fixed scaling $s$. Identical rewards ensure
performance differences are attributable to architecture.

\subsection{Training Algorithms and the Mechanics Ablation}
\label{sec:mechanics}
Each architecture is trained under \textbf{two} on-policy algorithms:
A2C~\cite{mnih2016a2c} and PPO~\cite{schulman2017ppo}, which differ
exactly in trust-region control (PPO's clipped surrogate bounds per-update
policy movement; A2C has no such bound). Both use GAE
($\lambda{=}0.95$)~\cite{schulman2016gae} with advantage normalisation,
entropy bonus, linear LR annealing ($1.0\times\rightarrow0.1\times$), and
Adam. Updates are computed on \emph{non-overlapping} batches of 20
transitions. Training runs 1000 epochs; evaluation is greedy (distribution
mean).

The non-overlapping-batch choice is itself a finding. An earlier version
of the pipeline used a stride-1 sliding window, producing an order of
magnitude more highly correlated overlapping updates per episode. Under
A2C this induced an early peak followed by sustained degradation
\emph{specific to Hard RA-RL} --- and, more insidiously, manufactured a
large, highly significant maximum-drawdown advantage for Soft RA-RL
($p{=}0.0025$ paired, $n{=}9$) that \emph{reversed sign} once the batching
was corrected, under both algorithms
(Section~\ref{sec:results-mechanics}). All results in this paper use the
corrected mechanics unless explicitly marked otherwise.

One residual asymmetry is disclosed: the L2 weight-penalty coefficient
was tuned per architecture during development (0.5 for Baseline/Hard,
0.01 for Soft/Router). \pending{L2-matched control run (Hard at 0.01)
queued; will either defuse or re-scope this limitation}

%% ===================================================================
\section{Experimental Setup}
%% ===================================================================

\subsection{Data and Universe}
38 large-cap U.S. equities (Appendix~\ref{app:universe}) with complete
daily history 2015-01-01 to 2024-01-01, from Yahoo Finance via
FinRL~\cite{liu2020finrl}; the fixed universe implies mild survivorship
bias, disclosed not corrected. The HMM basket is disjoint from the
trading universe.

\subsection{Protocols}
\textbf{Single-period:} train 2015--2021, test 2022-01-01 to 2024-01-01
(a bear year plus a recovery year); 3 seeds per (architecture, algorithm).
\textbf{Walk-forward:} 8 expanding-train windows, six-month non-overlapping
test periods 2020--2023 (COVID crash, recovery, 2022 bear, 2023 bull);
2 seeds per window; per-window HMM refit (no lookahead).

\subsection{Baselines}
Equal-weight (1/N), Markowitz mean--variance, S\&P~500 buy-and-hold, an
LSTM-context RL baseline (temporal context without an explicit regime
signal), and two \emph{HMM-gated rule} baselines consuming the same
signal as the RL agents: bear$\rightarrow$cash (hard rule) and equity
exposure $=1-P(\mathrm{Bear})$ (soft rule), otherwise 1/N.

\subsection{Statistics}
All architecture comparisons share seeds, so we report one-tailed
\emph{paired} $t$ and Wilcoxon signed-rank tests (pairing removes
between-seed variance; the unpaired Welch tests used in early versions of
this work are uniformly weaker). Walk-forward comparisons use Wilcoxon
over the 8 paired per-window Sharpes and Jobson--Korkie--Memmel
tests~\cite{jobson1981,memmel2003} on the pooled ${\sim}1000$-day daily
return series. \pending{Holm--Bonferroni across the metric family for the
headline table}

%% ===================================================================
\section{Results}
%% ===================================================================

\subsection{Single-Period: Soft vs.\ Hard under A2C}
\label{sec:results-single}

Table~\ref{tab:single-a2c} reports the corrected-mechanics single-period
comparison. Three patterns emerge. First, \textbf{soft gating dominates
in effect size}: Sharpe 0.462 vs 0.142, return $+26.2\%$ vs $+0.3\%$,
with soft winning two of three seed-pairs. At $n{=}3$ pairs these gaps
are not individually significant (paired $t$ $p_1{\approx}0.18$;
Wilcoxon's floor at $n{=}3$ is $0.125$); we report effect sizes and
per-seed behaviour rather than overclaiming.\footnote{Under the
\emph{original} mechanics this comparison was significant on all four
metrics at pooled $n{=}9$ --- but Section~\ref{sec:results-mechanics}
shows why we do not rest claims on those runs.} Second, \textbf{Hard
RA-RL is seed-unstable} (Figure~\ref{fig:seed-dots}): two of three seeds
collapse to negative Sharpe even under corrected mechanics, while every
soft seed is positive and
soft's dispersion is smaller on every ratio metric --- the instability
is a property of discrete switching under A2C, not an artifact of the
earlier update-scheme bug. Third, \textbf{the drawdown ordering is
reversed relative to the original mechanics} (Hard $48.9\%$ vs Soft
$57.2\%$, reversed in all three seed-pairs): see
Section~\ref{sec:results-mechanics}.

\begin{figure}[t]
\centering
%\includegraphics[width=\linewidth]{figures/seed_sharpe_dots.png}
\caption{Per-seed test Sharpe under both algorithms (3 shared seeds;
bars are means). Under A2C, two of three hard seeds go negative while
every soft seed is positive; under PPO both gated variants are rescued
and the learned router's collapse persists.
\pending{copy figures/seed\_sharpe\_dots.png into the project and
uncomment}}
\label{fig:seed-dots}
\end{figure}

\begin{table}[t]
\caption{Single-period test (2022--2024), corrected mechanics, A2C,
1000 epochs, $n{=}3$ shared seeds; one-tailed paired $t$ in the
soft-favoured direction (for max drawdown, soft-lower).}
\label{tab:single-a2c}
\begin{tabular}{lccc}
\toprule
Metric & Hard RA-RL & Soft RA-RL & paired $p_1$ \\
\midrule
Return (\%)        & $+0.3 \pm 28.0$   & $+26.2 \pm 29.8$  & $0.21$ \\
Sharpe             & $0.142 \pm 0.364$ & $0.462 \pm 0.245$ & $0.18$ \\
Max drawdown (\%)  & $48.9 \pm 3.9$    & $57.2 \pm 2.9$    & $0.92$ \\
Sortino            & $0.144 \pm 0.364$ & $0.469 \pm 0.242$ & $0.18$ \\
\bottomrule
\end{tabular}
\end{table}

\subsection{Learned Router: Collapse under Both Algorithms}
\label{sec:results-router}

Under A2C, the router produced large negative returns on two of three
seeds ($-19.0\%$, $-15.1\%$) and a negative mean Sharpe ($-0.060$),
below even the regime-agnostic baseline. Under PPO
(Table~\ref{tab:ppo}), the pattern \emph{replicates}: seed 2 collapses
($-27.8\%$, Sharpe $-0.281$), and the router's variance is
$1.6$--$2.5\times$ that of either gated variant on every metric. Ranking
Soft $>$ Hard $>$ Router holds under both algorithms. This is the
strongest cross-algorithm result in the paper: \textbf{jointly learning
the routing function destabilises training relative to consuming a
stable, separately estimated signal --- independently of the policy
optimiser}. \pending{extra PPO router/soft seeds ($n{=}3\rightarrow6$)
queued; current Soft-vs-Router $p_1\approx0.11$--$0.15$}

\subsection{Algorithm Robustness: PPO Localises Where Soft Gating Pays}
\label{sec:results-ppo}

\begin{table}[t]
\caption{Single-period test under PPO, $n{=}3$ seeds (mean $\pm$ std).
Welch one-tailed Soft-vs-Hard Sharpe $p_1{=}0.40$; pairing does not
change the conclusion ($p_1{=}0.42$).}
\label{tab:ppo}
\begin{tabular}{lccc}
\toprule
Metric & Hard & Soft & Router \\
\midrule
Return (\%)       & $40.7 \pm 16.7$ & $53.4 \pm 23.6$ & $14.3 \pm 38.8$ \\
Sharpe            & $0.726 \pm 0.214$ & $0.774 \pm 0.221$ & $0.314 \pm 0.538$ \\
Max drawdown (\%) & $36.4 \pm 4.5$  & $41.7 \pm 5.0$  & $41.7 \pm 11.9$ \\
Sortino           & $0.739 \pm 0.230$ & $0.794 \pm 0.236$ & $0.331 \pm 0.552$ \\
\bottomrule
\end{tabular}
\end{table}

Under PPO the Soft-vs-Hard Sharpe gap closes to a statistical tie
(0.774 vs 0.726) and Hard's seed instability disappears. We interpret
this mechanistically rather than as a failure to replicate: A2C applies
unconstrained on-policy updates, and discrete head-switching is
maximally exposed to them --- each head trains on roughly a quarter of
the signal density, and the trunk's gradient character switches abruptly
with the active head. PPO's clipped objective bounds per-update policy
movement and thereby \emph{rescues} discrete switching. The practical
design rule: \textbf{soft gating buys the most when the optimiser lacks
trust-region control; with PPO, the cheaper hard switch is competitive
on risk-adjusted return} (though the router remains broken under both).

\subsection{Walk-Forward Validation}
\label{sec:results-wf}

The walk-forward protocol delivers the study's most instructive
\emph{negative} result (Table~\ref{tab:wf}): with corrected mechanics,
the regime gate adds nothing. The no-gate baseline (aggregate Sharpe
1.270) statistically ties Hard (1.265) and edges Soft (1.105); no
RL-vs-RL comparison approaches significance under Wilcoxon over the 8
paired windows (all $p \geq 0.55$) or under Jobson--Korkie--Memmel on
the pooled ${\sim}1000$-day daily return series (e.g.\ Soft vs Baseline
$z{=}{-}0.60$, $p{=}0.55$). Equal weight remains the aggregate winner
(1.645, lowest drawdown) --- the only significant aggregate fact in the
protocol (vs.\ S\&P~500, $p{=}0.0078$).

Two window-level claims from the original (pre-correction) version of
this study die here, and we report them as cautionary results. First,
the COVID window: under the original mechanics, Soft appeared uniquely
crisis-robust (Sharpe 1.00 vs Hard 0.05 vs Baseline $-0.44$) ---
seemingly the textbook case for regime awareness. Under corrected
mechanics, \emph{every} variant weathers W1 (Baseline 1.11, Hard 1.12)
and Soft is the laggard (0.68). Second, Hard's anomalous H2-2022 spike
(2.83) normalises to 0.95. A transition-day decomposition ($\pm2$ days
around HMM state switches) shows all methods lose on transition days
(Sharpe $-0.9$ to $-2.0$, vs $+1.2$ to $+1.6$ on stable days), with
Hard nominally least-bad --- descriptive only, as between-method
transition differences are not significant.

\textbf{Synthesis with Section~\ref{sec:results-single}.} The two
protocols answer different questions, and together they localise when
regime-aware architecture pays. Trained \emph{once} on seven years and
asked to survive an unseen bear-plus-recovery (single-period), soft
gating shows a large advantage and hard switching is unstable. Retrained
\emph{every six months} on an expanding window (walk-forward), no gated
variant beats the no-gate baseline: \textbf{frequent retraining
substitutes for regime awareness}, because the policy's weights already
track the recent regime. Regime conditioning is a mechanism for
policies that must generalise across regime shifts they cannot be
retrained through.

\begin{table}[t]
\caption{Walk-forward aggregate (8 windows $\times$ 2 seeds, corrected
mechanics; mean across windows, seeds averaged within window).}
\label{tab:wf}
\begin{tabular}{lccc}
\toprule
Method & Sharpe & Return (\%) & MaxDD (\%) \\
\midrule
Equal-Weight        & \textbf{1.645} & 11.6 & \textbf{12.6} \\
Markowitz MVO       & 1.130 & 12.2 & 16.0 \\
S\&P~500 B\&H       & 1.104 & 6.4  & 13.7 \\
Baseline (no gate)  & 1.270 & 18.5 & 22.0 \\
Hard RA-RL          & 1.265 & 15.3 & 19.1 \\
Soft RA-RL          & 1.105 & 13.7 & 18.6 \\
\bottomrule
\end{tabular}
\end{table}

\subsection{Why RL at All? The Gated Equal-Weight Rule}
\label{sec:results-gatedew}

If the HMM signal were directly monetizable, a hand-crafted rule should
capture the value without RL. It does not (Table~\ref{tab:gatedew}): on
the single-period test, bear-gating \emph{cuts} equal-weight's Sharpe
from 0.41 to 0.16 (hard rule) or 0.19 (soft rule), while the soft-gated
RL agent's mean Sharpe (0.462) exceeds plain equal weight --- though the
latter comparison is seed-dependent (one of three seeds drives it) and
we claim it only at the mean level. On the walk-forward protocol the
rule is likewise (mildly) value-destroying. The asymmetry is the point:
the same signal that \emph{loses} money under the obvious fixed rule
adds value when the mapping from regime to allocation is
\emph{learned} --- which assets, how much, when.

\begin{table}[t]
\caption{HMM-gated equal-weight rules vs.\ plain 1/N (Sharpe).}
\label{tab:gatedew}
\begin{tabular}{lccc}
\toprule
Protocol & 1/N & Hard-gated & Soft-gated \\
\midrule
Walk-forward (agg.) & 1.645 & 1.559 & 1.561 \\
Single-period       & 0.410 & 0.161 & 0.188 \\
\midrule
Soft RA-RL (single-period) & \multicolumn{3}{c}{$0.462 \pm 0.245$ ($n{=}3$ seeds)} \\
\bottomrule
\end{tabular}
\end{table}

\subsection{The Mechanics Artifact}
\label{sec:results-mechanics}

Under the original stride-1 sliding-window updates, Soft RA-RL showed a
maximum-drawdown advantage over Hard that was the single most
significant result in the study (23.8\% vs 36.2\%, paired
$p{=}0.0025$, $n{=}9$). Under non-overlapping batches the ordering
\emph{reverses} under A2C (Hard $48.9\%$ vs Soft $57.2\%$, reversed in
all three seed-pairs) --- and independently reverses under PPO (36.4\%
vs 41.7\%), which never used the sliding buffer. The risk-adjusted-return
advantage, by contrast, survives the correction under A2C in direction
and effect size (Table~\ref{tab:single-a2c}).

The walk-forward protocol was hit even harder. Under the original
mechanics the no-gate baseline appeared broken (aggregate Sharpe 0.35),
Soft ranked best among RL variants (1.35), and Soft looked uniquely
crisis-robust in the COVID window --- three findings that together made
a clean pro-regime-gating story. Correcting the update scheme repaired
the \emph{baseline} most of all (0.35 $\rightarrow$ 1.27) and dissolved
all three claims (Section~\ref{sec:results-wf}). In total, four
publishable-looking results --- the drawdown advantage, the walk-forward
ranking, the COVID stress-test story, and Hard's H2-2022 anomaly ---
were artifacts of one update-scheme choice orthogonal to the
architectures under study. The lesson we draw is methodological:
\textbf{training-mechanics choices can manufacture an internally
consistent, statistically significant, but false architectural
narrative, and only a mechanics re-run catches it.} We recommend
mechanics ablations --- at minimum, buffer/batching variants --- as a
standard robustness pass for RL-in-finance comparisons, alongside seeds
and algorithm changes.

%% ===================================================================
\section{Discussion}
%% ===================================================================

\subsection{The 1/N Puzzle, Revisited}
Equal-weight achieves the highest aggregate walk-forward Sharpe of any
method, replicating DeMiguel et al.~\cite{demiguel2009}: five of eight
windows are bull or strong-recovery periods where full diversification
is near-optimal and active reallocation adds estimation noise. Our
claims are therefore scoped to \emph{active, RL-based} strategies. One
qualification: the 1/N dominance is protocol-dependent --- on the
2022--2024 single-period test, plain 1/N (Sharpe 0.410) is below the
soft-gated agent's mean (0.462), though seed-dependently. The
original-mechanics finding that Soft beat 1/N in the COVID window did
not survive the correction (Section~\ref{sec:results-wf}); in the
corrected walk-forward the RL variants beat 1/N in W1 (all ${\geq}0.68$
vs 0.33) but 1/N wins the aggregate.

\subsection{Interpretation: When Does a Regime Gate Pay?}
Four results triangulate one design principle. The router's
cross-algorithm collapse says regime \emph{detection} should be
external and stable, not jointly learned. The PPO tie says the
\emph{exploitation} mechanism (hard vs soft) matters mainly when the
optimiser is unconstrained. The gated-rule failure says exploitation
must be \emph{learned} --- the signal alone is not alpha. And the
walk-forward null says the gate only pays when the policy must
\emph{generalise across} regime shifts it cannot be retrained through:
with six-monthly retraining on expanding data, the network's weights
already encode the recent regime, and the gate is redundant. Together:
detect externally, exploit by learning, blend smoothly when training
lacks trust-region control --- and expect the benefit precisely in
deployment settings where retraining is infrequent relative to regime
turnover.

\subsection{Limitations}
(1)~Single asset class and fixed 38-stock universe with survivorship
bias; frictionless execution beyond a linear turnover penalty.
(2)~The L2 coefficient asymmetry (Section~\ref{sec:mechanics});
\pending{resolved by the L2-matched run if it completes}.
(3)~Walk-forward uses 2 seeds per window; Wilcoxon at $n{=}8$ windows
has a $p$-floor of 0.0078.
(4)~The evaluation window (2020--2024) contains exactly one
flight-to-quality crash and one rate-driven bear; a probe of fixed
defensive-rotation rules shows they help in the former (+0.30
window Sharpe, COVID) and destroy value in the latter ($-1.21$, 2022,
JK--Memmel $p{=}0.042$ against the rule) --- regime \emph{type}
matters, motivating multi-asset extensions with richer regime
detection as future work.
(5)~Absolute Sharpe levels are modest by institutional standards; the
comparison is relative, under an identical reward.

%% ===================================================================
\section{Conclusion}
%% ===================================================================

Given an external HMM regime signal, we compared hard switching, soft
posterior gating, and end-to-end learned routing for RL portfolio
policies, under two optimisers, two evaluation protocols, seed-paired
statistics, and a training-mechanics ablation. The learned router is
the worst and least stable variant under both A2C and PPO --- anchor
the gate to a stable, separately estimated signal. Soft gating's
advantage over hard switching is real but doubly conditional: large
under unconstrained A2C updates yet a tie under PPO's clipped
objective, and present only when the policy is trained once and must
generalise across unseen regime shifts --- under six-monthly
walk-forward retraining, no gated variant beats a no-gate baseline,
because frequent retraining substitutes for regime awareness.
Hand-crafted rules on the same signal destroy value, so the signal pays
only through learned exploitation. And an entire pro-gating narrative
of our own early results --- a significant drawdown advantage, a
walk-forward ranking, a COVID stress-test story --- proved to be
artifacts of a sliding-window update buffer, dissolving under corrected
mechanics. We offer the full protocol (paired seeds, two algorithms,
two protocols, mechanics ablation, rule baseline) as a template for
architectural claims in RL for finance. Future work: multi-asset
regime rotation with bear-type-aware detection, differential-Sharpe
rewards~\cite{moody2001}, and off-policy algorithms (a SAC suite is
implemented and calibrated; its full evaluation is deferred).

%% ===================================================================
%% Acknowledgments — SUPPRESSED for double-blind review. Restore for
%% camera-ready: URECA programme funding + supervisor thanks.
%% ===================================================================

\bibliographystyle{ACM-Reference-Format}
\bibliography{references}

\appendix

\section{Walk-Forward Windows}
\label{app:windows}
Eight expanding-train windows, all starting 2015-01-01; test halves
H1-2020 (COVID crash), H2-2020 (recovery), H1/H2-2021 (bull),
H1-2022 (bear onset), H2-2022 (bear/recovery), H1/H2-2023 (bull).

\section{Trading Universe}
\label{app:universe}
AAPL ABBV ABT AMD AMZN BA BAC C CAT COST CRM CVX DIS GE GOOGL GS HD HON
INTC JNJ JPM KO MCD META MMM MRK MS MSFT NFLX NVDA PEP PFE PG TXN UNH V
WMT XOM.

\section{Hyperparameters}
\label{app:hyper}
Adam; LR $10^{-4}$ (Baseline/Hard) and $3\times10^{-5}$ (Soft/Router),
linearly annealed to $0.1\times$; $\gamma{=}0.99$; GAE $\lambda{=}0.95$
with advantage normalisation; non-overlapping batch 20; value-loss
coefficient 0.5; entropy 0.01; grad-clip 0.5; L2 penalty 0.5
(Baseline/Hard) / 0.01 (Soft/Router); trunk $2\times256$ ReLU; 1000
epochs. PPO: clip $\epsilon{=}0.2$, 10 optimisation epochs per rollout,
minibatch 64 over full-episode rollouts, LR $10^{-4}$ (all variants),
same $\gamma$/GAE/value/entropy coefficients and LR annealing as A2C.
Reward: $\lambda{=}0.5$, $c{=}10^{-4}$, $\kappa{=}0.005$, bounds
$[-0.05,0.20]$.

\section{Per-Seed Results}
\label{app:perseed}

\begin{table}[h]
\caption{Per-seed results, corrected-mechanics A2C single-period run
(underlying Table~\ref{tab:single-a2c}).}
\label{tab:perseed-a2c}
\begin{tabular}{llcccc}
\toprule
Variant & Seed & Return (\%) & Sharpe & MaxDD (\%) & Sortino \\
\midrule
Hard & 0 & $-17.50$ & $-0.083$ & $50.72$ & $-0.080$ \\
Hard & 1 & $-14.21$ & $-0.053$ & $51.54$ & $-0.053$ \\
Hard & 2 & $+32.56$ & $+0.563$ & $44.36$ & $+0.565$ \\
Soft & 0 & $+3.01$  & $+0.262$ & $57.17$ & $+0.267$ \\
Soft & 1 & $+59.88$ & $+0.735$ & $54.40$ & $+0.737$ \\
Soft & 2 & $+15.82$ & $+0.390$ & $60.12$ & $+0.403$ \\
\bottomrule
\end{tabular}
\end{table}

\begin{table}[h]
\caption{Per-seed results, PPO single-period run (underlying
Table~\ref{tab:ppo}).}
\label{tab:perseed-ppo}
\begin{tabular}{llcccc}
\toprule
Variant & Seed & Return (\%) & Sharpe & MaxDD (\%) & Sortino \\
\midrule
Hard   & 0 & $+51.36$ & $+0.886$ & $31.24$ & $+0.918$ \\
Hard   & 1 & $+49.22$ & $+0.809$ & $39.59$ & $+0.819$ \\
Hard   & 2 & $+21.44$ & $+0.483$ & $38.35$ & $+0.479$ \\
Soft   & 0 & $+28.07$ & $+0.538$ & $46.99$ & $+0.534$ \\
Soft   & 1 & $+74.86$ & $+0.975$ & $40.92$ & $+0.995$ \\
Soft   & 2 & $+57.36$ & $+0.808$ & $37.10$ & $+0.854$ \\
Router & 0 & $+21.94$ & $+0.454$ & $36.40$ & $+0.461$ \\
Router & 1 & $+48.66$ & $+0.767$ & $33.46$ & $+0.807$ \\
Router & 2 & $-27.80$ & $-0.281$ & $55.35$ & $-0.274$ \\
\bottomrule
\end{tabular}
\end{table}

Under A2C (original-mechanics campaign, 1000 epochs, shared seeds), the
Learned Router's three seeds returned $-18.96\%$, $+18.80\%$, and
$-15.05\%$ (Sharpe $-0.398$, $+0.458$, $-0.241$) --- two of three
collapsed, mirroring the PPO seed-2 collapse above.

%% Figures to include (all generated in Final/figures/; corrected-mechanics
%% versions come from make_icaif_plots.py):
%%   fig:seed-dots    — seed_sharpe_dots.png (CORRECTED; the narrative figure)
%%   fig:equity       — equity_curve_corrected_seed1.png (CORRECTED)
%%   fig:bic          — hmm_state_selection BIC/AIC sweep
%%   fig:gate         — gate_posterior.png (+ COVID zoom)
%%   fig:architectures— hard/soft/router diagrams (compress to one row)
%%   fig:wf-sharpe    — per-window Sharpe (REGENERATE once batchfix wf lands;
%%                      walkforward_sharpe.png currently plots superseded data)
%%   bar_corrected_a2c.png also available if a bar form of Table 1 is wanted
\begin{figure}[t]
\centering
%\includegraphics[width=\linewidth]{figures/hmm_state_selection.png}
\caption{HMM model selection: BIC favours $K{=}4$; $K\geq5$ improves BIC
only via degenerate states occupying $0.3$--$2\%$ of days.
\pending{copy figure from Final/figures/ and uncomment}}
\label{fig:bic}
\end{figure}

\end{document}
